\section{Objectif du projet}

% ----------------------------------------------------------------
% 📌 À quoi sert cette section ?
% ----------------------------------------------------------------
% Cette section doit répondre clairement à la question :
% « Que souhaite-t-on accomplir avec ce projet ? »
% 
% Il ne s’agit pas encore de détailler les fonctionnalités ni les solutions techniques, 
% mais de donner une vision synthétique et globale des finalités du projet.
%
% L’objectif doit être SMART (Spécifique, Mesurable, Atteignable, Réaliste, Temporellement défini),
% ou du moins respecter ces grandes lignes pour aider le lecteur à comprendre l’utilité du projet.

% ----------------------------------------------------------------
% 🧭 Comment la remplir ?
% ----------------------------------------------------------------
% ➤ Reformulez l’objectif du projet en une ou deux phrases d’introduction.
% ➤ Ajoutez ensuite un développement précisant :
%     - Le public cible
%     - Le contexte d’utilisation
%     - Le besoin ou la motivation ayant mené à ce projet
%     - Les grandes lignes du résultat attendu
%
% ➤ Vous pouvez, si cela aide à clarifier, ajouter un schéma illustrant le positionnement
%   ou la finalité du projet dans son environnement.

% ----------------------------------------------------------------
% 🧾 Exemple rédigé :
% ----------------------------------------------------------------
% L’objectif du projet est de concevoir un outil interactif permettant aux étudiants de s’initier
% au langage SVG à travers une interface ludique de type "bac à sable".
% Le but est de rendre l’apprentissage de la syntaxe SVG plus intuitif, tout en conservant
% une liberté de création.
%
% Ce projet s’adresse aux étudiants de première année en informatique, dans un contexte éducatif.
% Le livrable visé est une plateforme accessible en ligne, ne nécessitant aucune installation locale,
% et intégrant des défis progressifs ainsi qu’une prévisualisation en temps réel.

% ----------------------------------------------------------------
% 📊 Exemple de schéma à insérer (schéma positionnement)
% ----------------------------------------------------------------

% \begin{figure}[H]
%     \centering
%     \begin{tikzpicture}[node distance=1.5cm]
%         \node (start) [draw, rectangle] {Problème identifié : apprentissage SVG difficile};
%         \node (middle) [below of=start, draw, rectangle] {Solution proposée : plateforme interactive};
%         \node (end) [below of=middle, draw, rectangle] {Objectif final : apprentissage ludique et progressif};
%         \draw[->] (start) -- (middle);
%         \draw[->] (middle) -- (end);
%     \end{tikzpicture}
%     \caption{Schéma simplifié de l’objectif du projet}
%     \label{fig:objectif}
% \end{figure}

% ----------------------------------------------------------------
% ✍️ À vous de jouer : complétez ci-dessous
% ----------------------------------------------------------------

% (Exemple de rédaction)
L’objectif du projet est de...

